\SourceLeaf{32}\noindent\hspace*{\parindent}Die unter [1] als »Theorem I« angegebenen Formeln sind hiernach leicht
ersichtlich; und es ist freilich nicht gewiss, aber sehr wahrscheinlich, dass
\textsc{Gauss} dieselben unter Benutzung der complexen Factoren $t_1$, $t_2$
aufgestellt hat.

Welche Absicht \textsc{Gauss} mit der Bildung der Zahlenreihe $t'$, $t''$, $t'''$,
.\,.\,. gehabt hat, ist nicht angegeben.

In den weiter folgenden Entwicklungen sind an Stelle von $t$, $\varphi t$,
$Ft$ die Bezeichnungen $\Delta$, $D$, $\nu\nu A+\nu B+C$ getreten. Es handelt
sich darum, die gegenseitige Beziehung der beiden Zahlen $a+b\nu+c\nu\nu$ und
$\nu\nu A+\nu B+C$ weiter zu verfolgen. Nach handschriftlichen Aufzeichnungen
\textsc{Scherings} kann man diese Formeln unmittelbar aus den Gleichungen des
»Theorems I« ablesen. Setzt man z.\,B. in

\[
F(Ft) = (\varphi t)\,t
\]

\noindent die jetzige Bedeutung von $t$, $\varphi t$ und $Ft$ ein, so folgt:

\[
F(Ft) = (BB-AC)\nu\nu + (nAA-BC)\nu + (CC-nAB)
\]
\[
(\varphi t)\,t = D(c\nu\nu+b\nu+a),
\]

\noindent sowie aus der Identität beider Ausdrücke:

\[
CC-nAB = Da, \qquad nAA-BC = Db, \qquad BB-AC = Dc.
\]

In \textsc{Gauss}' Handschrift findet sich bei der Notiz [2] noch die Angabe:

\begin{center}\textit{Correction von $\nu$ à peu près} $= -\dfrac{D^2}{9BCC},$\end{center}

\noindent sowie der unvollendete Satz:

\begin{center}\textit{Forma ternaria per subst.}\end{center}
\[
\begin{aligned}
x &= ax'+by'+cz', \\
y &= ncx'+ay'+bz', \\
z &= nbx'+ncy'+az',
\end{aligned}
\]

\Needspace{9\baselineskip}\noindent welcher dahin zu deuten ist, dass die ternäre Form

\[
nnx^3+ny^3+z^3-3nxyz
\]

\noindent vermöge jener Substitution übergeht in

\[
D(nnx'^3+ny'^3+z'^3-3nx'y'z').
\]

Der Satz [3] über cubische Reste entspringt direct aus den durch die
voraufgehenden Formeln begründeten Congruenzen:

\[
nA^3-B^3 \equiv 0, \qquad nB^3-C^3 \equiv 0, \qquad C^3-nnA^3 \equiv 0 (\mathrm{mod.}\ D).
\]

Endlich liefern die beiden numerischen Tabellen am Schlusse der Fragmente,
abgesehen von der für $n=8$ angegebenen Zahl $-2+\sqrt[3]{64}$ und den beiden
für $n=7$ und $n=9$ gemachten Angaben, bei denen indess nur Schreibfehler
vorzuliegen scheinen, lauter Zahlen $a+b\nu+c\nu\nu$ von der Norm $\pm 1$,
d.\,i. lauter Einheiten. In der ersten Tabelle, derjenigen für $n=2$, ist
$1+\sqrt[3]{2}+\sqrt[3]{4}$ die Fundamentaleinheit des Körpers, deren erste
neun Potenzen sammt ihren reciproken Einheiten berechnet werden. Einer
brieflichen%
